\subsection*{Part E: interpretative limitations}
The manuscript's one-word margin ``bihexagones'' is not accompanied by a
definition in the supplied pages. E.003 gives and proves a complete
explicit construction using pairs of edge-disjoint Hamiltonian hexagons.
It does not claim to determine uniquely what Grothendieck intended by
that word. Likewise E.010 states exact category conventions for both
``isotopique'' and ``isotypique'' readings; raw homeomorphism groupoids
are distinguished from their isotopy quotients.

Official scans 119/156,157,159,160,161 were opened and visually
inspected on 2026-11-19. The boxed formula on 157 explicitly has the
side-dual marker $Q^\vee$, confirming E.006. The normalizer subscript
in the cramped margin of 156 remains palaeographically uncertain;
E.008 proves both readings separately, so this uncertainty does not
leave an algebraic assertion unresolved.
